\section*{Chapter \ref{ch:b3:complex-mechanisms} --- Reaction Mechanisms of Complexes}

\begin{solution}{exo:b3:complex-mechanisms:1}
Inert: $\ce{[Cr(H2O)6]^{3+}}$ ($d^3$) and $\ce{[Co(NH3)6]^{3+}}$ (low-spin $d^6$), large ligand-field
activation energies. \omterm{def:b3:complex-mechanisms:labile}{Labile}: $\ce{[Cr(H2O)6]^{2+}}$ (high-spin $d^4$) and $\ce{[Cu(H2O)6]^{2+}}$ ($d^9$),
Jahn--Teller distorted with long, weak axial bonds; $\ce{[Zn(H2O)6]^{2+}}$ ($d^{10}$), no ligand-field
barrier.
\end{solution}

\begin{solution}{exo:b3:complex-mechanisms:2}
When $k_2[\mathrm Y] \gg k_{-1}[\mathrm X]$ the denominator is $k_2[\mathrm Y]$ and $v = k_1[\mathrm{ML_5X}]$; when $k_{-1}[\mathrm X] \gg k_2[\mathrm Y]$ it is
$k_{-1}[\mathrm X]$ and the second form follows (a pre-equilibrium inhibited by the leaving
ligand).
\end{solution}

\begin{solution}{exo:b3:complex-mechanisms:3}
D: $\Delta V^{\ddagger} > 0$ and $\Delta^{\ddagger}S^\circ > 0$ (a ligand leaves). A: both negative (a ligand is bound in the
transition state).
\end{solution}

\begin{solution}{exo:b3:complex-mechanisms:4}
$cis$-$\ce{[PtCl2(NH3)2]}$ from $\ce{[PtCl4]^{2-}}$; $trans$-$\ce{[PtCl2(NH3)2]}$ from $\ce{[Pt(NH3)4]^{2+}}$ (the
\omterm{def:b3:complex-mechanisms:trans-effect}{trans effect} of \ce{Cl-} exceeds that of \ce{NH3}).
\end{solution}

\begin{solution}{exo:b3:complex-mechanisms:5}
$k_{\mathrm{obs}} = 2.0 \times \num{3.0e4} \times 0.10/1.2 = \qty{5.0e3}{s^{-1}}$; at \qty{1.0}{mol/L}, $\num{6.0e4}/3.0 = \qty{2.0e4}{s^{-1}}$; limit
$k_i = \qty{3.0e4}{s^{-1}}$.
\end{solution}

\begin{solution}{exo:b3:complex-mechanisms:6}
$\Delta V^{\ddagger} = -RT\ln2/\Delta p = -8.314 \times 298 \times 0.693/\num{99.9e6} = \qty{-1.7e-5}{m^3.mol^{-1}} = \qty{-17}{cm^3.mol^{-1}}$:
associative.
\end{solution}

\begin{solution}{exo:b3:complex-mechanisms:7}
The phosphine, a strong $\sigma$ donor and $\pi$ acceptor, weakens the bond trans to it (\omterm{def:b3:complex-mechanisms:trans-effect}{trans
influence}, longer bond) and labilises it (\omterm{def:b3:complex-mechanisms:trans-effect}{trans effect}): that chloride is
replaced first.
\end{solution}

\begin{solution}{exo:b3:complex-mechanisms:8}
The bridging ligand ends up transferred to the oxidised metal; the rate depends
strongly on the nature of the potential bridge; a binuclear intermediate is
detected; and the rate cannot exceed that of the substitution needed to form the
bridge on the \omterm{def:b3:complex-mechanisms:labile}{labile} partner.
\end{solution}

\begin{solution}{exo:b3:complex-mechanisms:9}
$(\lambda + \Delta G^\circ)^2/4\lambda$ with $4\lambda = 320$: 20.0, 11.3, 0 and \qty{5.0}{kJ/mol} (the last in the \omterm{def:b3:complex-mechanisms:inverted}{inverted
region}).
\end{solution}

\begin{solution}{exo:b3:complex-mechanisms:10}
$k_{12} = \sqrt{4.0 \times \num{2.0e5} \times \num{1.0e4}} = \qty{8.9e4}{L.mol^{-1}.s^{-1}}$.
\end{solution}

\begin{solution}{exo:b3:complex-mechanisms:11}
High-spin $d^5$: 0; $d^7$: $1.33\,Dq$; $d^8$: $2.0\,Dq$. Expected water-exchange rates: \ce{Mn^{2+}} $>$ \ce{Co^{2+}} $>$
\ce{Ni^{2+}}, the order observed.
\end{solution}

\begin{solution}{exo:b3:complex-mechanisms:12}
For very exergonic reactions the intrinsic rate exceeds the rate at which the ions
meet: the observed constant stays at the diffusion limit and the decrease is
hidden. With donor and acceptor fixed at a set distance in one molecule, there
is no diffusion step, and a series of acceptors of increasing driving force
showed the rate rising, passing a maximum and falling.
\end{solution}

\begin{solution}{pb:b3:complex-mechanisms:1}
\textbf{1.} $5d^8$.
\textbf{2.} The strong field of a $5d$ ion makes the $d_{x^2-y^2}$ orbital, pointing at four
ligands, very high: eight electrons pair in the four lower orbitals.
\textbf{3.} The complex has 16 valence electrons and an open axial site: the entering
ligand binds first, through a five-coordinate transition state.
\textbf{4.} $k_{\mathrm{obs}} = k_1 + k_2[\mathrm Y]$: a solvent path (solvent attack, then fast replacement by Y)
and a direct path.
\textbf{5.} Substitution takes hours (part III): inert on the scale of minutes, not of
a day.
\textbf{6.} Cisplatin, because the second \ce{NH3} replaces a chloride trans to chloride
(\omterm{def:b3:complex-mechanisms:trans-effect}{trans effect} $\ce{Cl-} > \ce{NH3}$); side products form as well.
\textbf{7.} Iodide has a larger \omterm{def:b3:complex-mechanisms:trans-effect}{trans effect} than chloride, which makes the directing
effect cleaner and faster.
\textbf{8.} In $\ce{[PtI3(NH3)]-}$ the iodide trans to iodide is labilised, not the one trans to
\ce{NH3}: the second ammonia enters cis to the first.
\textbf{9.} Silver iodide precipitates and $cis$-$\ce{[Pt(H2O)2(NH3)2]^{2+}}$ remains in solution.
\textbf{10.} Chloride replaces the two water molecules at their own positions; the
ammines are not touched.
\textbf{11.} From $\ce{[Pt(NH3)4]^{2+}}$ and two chlorides.
\textbf{12.} Its two reactive positions are opposite each other: it cannot bind two
neighbouring bases of the same DNA strand, the lesion through which cisplatin
acts.
\textbf{13.} $\ln2/\num{9.0e-5} = \qty{7.7e3}{s} = \qty{2.1}{h}$.
\textbf{14.} $1 - \eu^{-\num{9.0e-5} \times 7200} = 0.48$.
\textbf{15.} Water is a much better leaving group than chloride: the aqua complex is
replaced quickly by the nitrogen of a DNA base.
\textbf{16.} Associative, as for all platinum(II) substitutions: negative activation
volume.
\textbf{17.} The high chloride concentration pushes the equilibrium back towards the
dichloro form, keeping the drug intact until it reaches the cells.
\textbf{18.} $f = K/(K + [\ce{Cl-}])$.
\textbf{19.} $\num{3.0e-3}/0.103 = 0.029$.
\textbf{20.} $\num{3.0e-3}/0.0070 = 0.43$.
\textbf{21.} 15.
\textbf{22.} Only inside cells does an appreciable fraction exist as the reactive aqua
complex.
\textbf{23.} Sulfur ligands, such as glutathione and the cysteine and methionine of
proteins, which bind the soft platinum strongly.
\textbf{24.} Ammonia binds platinum strongly and lies trans to chloride, a ligand of
small \omterm{def:b3:complex-mechanisms:trans-effect}{trans effect}: the ammine bonds are not labilised.
\textbf{25.} \textbf{About 15 times more of the drug is in the reactive aqua form inside a
cell than in plasma.}
\end{solution}
